\documentclass[11pt,a4paper]{article}
\usepackage[T1]{fontenc}
\usepackage[utf8]{inputenc}
\usepackage{lmodern}
\usepackage[margin=29mm,headheight=14pt]{geometry}
\usepackage{amsmath,amssymb,amsthm,mathtools}
\usepackage{microtype,booktabs,longtable,enumitem,array}
\usepackage[hidelinks,breaklinks=true]{hyperref}
\usepackage{url}
\hypersetup{pdftitle={Two routes from a zero-cancellation condition to the Riemann hypothesis},pdfsubject={Conditional equivalences and their Lean formalization}}
\setlength{\emergencystretch}{2em}
\setlength{\parindent}{1.2em}
\setlength{\parskip}{2pt}
\setlist{itemsep=3pt,topsep=5pt}
\newtheorem{theorem}{Theorem}[section]
\newtheorem{lemma}[theorem]{Lemma}
\newtheorem{proposition}[theorem]{Proposition}
\newtheorem{corollary}[theorem]{Corollary}
\theoremstyle{definition}
\newtheorem{definition}[theorem]{Definition}
\newtheorem{remark}[theorem]{Remark}
\numberwithin{equation}{section}
\newcommand{\RH}{\mathrm{RH}}
\newcommand{\Ree}{\operatorname{Re}}
\newcommand{\Imm}{\operatorname{Im}}
\newcommand{\code}[1]{\texttt{\small #1}}
\title{\bfseries Two routes from a zero-cancellation condition\\
to the Riemann hypothesis}
\author{}
\date{Preprint\\28 October 2026}
\begin{document}
\maketitle
\vspace{-1.2em}
\begin{abstract}
We study a zero-cancellation condition for the Bernstein rows associated
with the actual normalized logarithmic-derivative moments of the Riemann
$\xi$-function. The condition, denoted $H_{00}$, asserts admissibility of
the cutoff-zero, coefficient-zero cancellation budget. We present two
constructions of the implication $H_{00}\Rightarrow\RH$. The first
passes from constant row variation to an existing actual spectral
criterion. The second constructs a Hausdorff probability measure, identifies
its resolvent with the actual logarithmic derivative near the origin,
and continues an analytic defect to exclude nonreal zeros. A retained
converse gives $H_{00}\Longleftrightarrow\RH$ for each construction.
The accompanying Lean development exports both implications with no
additional mathematical hypothesis. The condition $H_{00}$ remains a
universally quantified assertion: neither its unconditional validity nor
a terminating finite certificate checker is established here.
\end{abstract}
\vspace{0.5em}
\noindent\textbf{Keywords.} Riemann hypothesis; Hausdorff moments; Bernstein
weights; logarithmic derivative; positive measure; formal proof.

\section{Introduction}
Moment positivity provides a natural meeting point between finite
algebraic identities and the analytic location of zeros. In the present
setting, the finite objects are Bernstein rows formed from a fixed,
actual sequence of normalized logarithmic-derivative moments. Their
total variation is at least one. Equality at every order forces all row
weights to be nonnegative and is equivalent to the vanishing of a
nonnegative cancellation quantity.

This note records two ways to pass from that cancellation condition to
the Riemann hypothesis. They share the actual moment sequence and the
finite positivity argument, but use different forward analytic
constructions. The first applies an already formalized spectral
characterization of subexponential row variation. The second constructs
a positive measure directly from the Bernstein rows and uses its
resolvent to control the logarithmic derivative.

The distinction between a finite description and a finite verification
is central. The candidate budget consists of two natural numbers, namely
$(0,0)$, but its admissibility quantifies over every row index. The
equivalence proved here therefore does not supply an accepted finite
certificate of RH. In particular, no proof of $H_{00}$ is claimed.

The mathematical development below is accompanied by the Lean sources
in~\cite{source}. We explicitly identify the previously formalized
actual-function and spectral results used as inputs. The note does not
reprove the full analytic infrastructure of that source development.
Section~\ref{sec:formal} distinguishes the mathematical statements,
reported native verification, and the separately configured GitHub CI
target.

\section{Actual analytic data}\label{sec:actual}
Let $\xi$ be the actual entire Riemann $\xi$-function used in the
accompanying development. Away from $s=0,1$, its normalization is
\begin{equation}\label{eq:xi}
 \xi(s)=\frac{s(s-1)}2\,\Lambda(s),
\end{equation}
where $\Lambda$ denotes the completed Riemann zeta function.
More exactly, with Mathlib's entire pole-subtracted auxiliary function
$\Lambda_0$ (\code{completedRiemannZeta}$_0$), the formal definition is
\[
 \xi(s)=\frac{1+s(s-1)\Lambda_0(s)}2.
\]
This specifies the values at the removable singularities as well. We use the retained equivalence between
Mathlib's literal \code{RiemannHypothesis} and the assertion that every
zero of this $\xi$ has real part $1/2$.

The reflection identity $\xi(1-s)=\xi(s)$ makes
$w\mapsto\xi(1/2+w)$ even. The actual function
\[
 \Theta(z)=\xi\bigl(\tfrac12+\sqrt z\bigr)
\]
is proved entire in the source, despite the square root used in its
definition, and satisfies
\begin{equation}\label{eq:centered-square}
 \Theta(w^2)=\xi\bigl(\tfrac12+w\bigr)\qquad(w\in\mathbb C).
\end{equation}
The following actual-function facts are retained from that development:
\begin{enumerate}[label=(\roman*)]
 \item $\Theta(0)\ne0$, and $\Theta$ is entire;
 \item the zero set of $\Theta$ is invariant under complex conjugation;
 \item if $\xi(s)=0$, then $|\Imm s|>1$;
 \item the moments defined below are real before taking real parts, and
 their normalizing constant is strictly positive.
\end{enumerate}
These are supplied lemmas about the actual functions, not additional
hypotheses of the final $H_{00}$ implications. Their source locations are
listed in Appendix~\ref{app:source}. Only the nonvanishing of $\Imm s$,
rather than the numerical strength of (iii), is needed in the final
return from $\Theta$-zeros to $\xi$-zeros.

On a neighborhood of zero define
\begin{equation}\label{eq:moments}
 R(z)=\frac{\Theta'(z)}{\Theta(z)},\qquad
 m_n=\Ree\!\left(\frac{(-1)^n}{n!}R^{(n)}(0)\right),\qquad
 \kappa=m_0>0,\qquad u_n=\frac{m_n}{\kappa^{n+1}}.
\end{equation}
Here $R^{(0)}=R$. In particular $u_0=1$. The retained reality result
turns~\eqref{eq:moments} into the complex derivative identity
\begin{equation}\label{eq:actual-derivatives}
 R^{(n)}(0)=(-1)^n n!\,\kappa^{n+1}u_n.
\end{equation}
Throughout the paper, unless a general real sequence is explicitly
introduced, $u$ denotes this actual normalized sequence.

\section{Bernstein cancellation and the main statement}\label{sec:cancel}
For a real sequence $u$, define its Bernstein weights by
\begin{equation}\label{eq:weights}
 w_{N,j}=\binom Nj\sum_{r=0}^{N-j}(-1)^r
                 \binom{N-j}{r}u_{j+r}\qquad(0\le j\le N),
 \qquad V_N=\sum_{j=0}^N|w_{N,j}|.
\end{equation}
The elementary row-mass and coarsening identities are
\begin{align}
 \sum_{j=0}^N w_{N,j}&=u_0,\label{eq:row-mass}\\
 (n+1)w_{n,j}
   &=(n+1-j)w_{n+1,j}+(j+1)w_{n+1,j+1}
       \quad(0\le j\le n).\label{eq:coarsening}
\end{align}
For $0\le j\le n$ put
$a_{n,j}=(n+1-j)w_{n+1,j}$ and
$b_{n,j}=(j+1)w_{n+1,j+1}$. Define the row cancellation by
\begin{equation}\label{eq:cancellation}
 C_n=\sum_{j=0}^n
 \bigl(|a_{n,j}|+|b_{n,j}|-|a_{n,j}+b_{n,j}|\bigr).
\end{equation}
The triangle inequality gives $C_n\ge0$. Summing the weighted terms and
using~\eqref{eq:coarsening} yields
\begin{equation}\label{eq:variation-increment}
 C_n=(n+1)(V_{n+1}-V_n).
\end{equation}

\begin{definition}[The semantic zero budget]
A cancellation budget is a pair $W=(N_0,b)\in\mathbb N^2$.
For the actual sequence it is \emph{admissible} if
\[
 \forall n\ge N_0,\qquad C_n\le bV_n.
\]
The condition $H_{00}$ is admissibility of $W=(0,0)$.
\end{definition}

\begin{proposition}\label{prop:finite-equivalences}
For the actual normalized sequence, the following are equivalent:
\begin{enumerate}[label=(\roman*)]
 \item $H_{00}$;
 \item $C_n=0$ for every $n\ge0$;
 \item $V_N=1$ for every $N\ge0$;
 \item $w_{N,j}\ge0$ for every $0\le j\le N$.
\end{enumerate}
\end{proposition}
\begin{proof}
Since $C_n\ge0$, admissibility of $(0,0)$ is equivalent to $C_n=0$
at every order. Equation~\eqref{eq:variation-increment} then says
$V_{n+1}=V_n$. Since $V_0=|u_0|=1$, induction gives (iii).
Conversely, (iii) forces the increment to vanish and hence gives (ii).

By~\eqref{eq:row-mass}, each row has mass one. Therefore
\[
 V_N-1=\sum_{j=0}^N\bigl(|w_{N,j}|-w_{N,j}\bigr).
\]
Every summand is nonnegative, and it vanishes precisely when the
corresponding weight is nonnegative. This proves the equivalence of
(iii) and (iv).
\end{proof}

\begin{theorem}[Two forward constructions and a common converse]
\label{thm:main}
For the actual sequence~\eqref{eq:moments},
\[
 H_{00}\quad\Longleftrightarrow\quad\RH.
\]
The forward implication admits both the spectral construction of
Section~\ref{sec:spectral} and the measure--resolvent construction of
Sections~\ref{sec:measure}--\ref{sec:closure}. Each construction has
$H_{00}$ as its only mathematical hypothesis. The two equivalences use
the same retained converse.
\end{theorem}
The proof is completed below. The theorem is an equivalence of
mathematical assertions. It does not assert that either assertion holds
unconditionally.

\section{The spectral route and the converse}\label{sec:spectral}
Call a nonnegative sequence $V$ \emph{subexponential} if
\begin{equation}\label{eq:subexponential}
 \forall\varepsilon>0\ \exists A>0\ \forall n\ge0,\qquad
 V_n\le A e^{\varepsilon n}.
\end{equation}
The earlier actual spectral development supplies the following
characterization for the specific $V$ in~\eqref{eq:weights}:
\begin{equation}\label{eq:retained-spectral}
 V\text{ is subexponential}\quad\Longleftrightarrow\quad\RH.
\end{equation}
This characterization is supplied by the actual spectral theorem
in~\cite{source}; the exact declaration is listed in
Appendix~\ref{app:source}. Its analytic inputs are constructed in that
development rather than assumed at the present endpoints.

More precisely, the canonical actual divisor gives complex nodes
$(q_i)_{i\in J}$, with zero multiplicities retained, satisfying
\begin{equation}\label{eq:spectral-calibration}
 q_i\ne0,\qquad \sum_i|q_i|<\infty,\qquad
 u_n=\sum_i q_i^{n+1},\qquad
 \alpha_i=-\frac{1}{\kappa q_i}.
\end{equation}
Every $\alpha_i$ is a zero of $\Theta$, and every zero occurs in this
family. Put
\[
 \rho=\sup\bigl(\{1\}\cup\{|q_i|+|1-q_i|:i\in J\}\bigr).
\]
The retained extraction theorem, applied to the constructed calibration,
proves the ordinary limit
\begin{equation}\label{eq:spectral-limit}
 \lim_{N\to\infty}V_N^{1/N}=\rho.
\end{equation}
Its proof separates a finite dominant exponential part from a tail
with strictly smaller growth; the full extraction argument is in the
accompanying source. Since $|q|+|1-q|\ge1$, with equality precisely for
$q\in[0,1]$, the case $\rho=1$ is equivalent to all nodes lying in that
segment. The actual calibration identifies this case with RH and yields
~\eqref{eq:retained-spectral}. We use that already constructed spectral
result here, without reproducing its full proof.

\begin{proof}[Proof of the first forward implication]
Assume $H_{00}$. Proposition~\ref{prop:finite-equivalences} gives
$V_n=1$ for every $n$. For every $\varepsilon>0$, take $A=1$ in
~\eqref{eq:subexponential}, since
$1\le e^{\varepsilon n}$ for $n\ge0$. The forward direction of
~\eqref{eq:retained-spectral} proves $\RH$.
\end{proof}

For the common converse, the actual calibration under RH has every
$q_i\in[0,1]$. Substitution into the finite Bernstein formula gives
\[
 w_{N,j}=\binom Nj\sum_i q_i^{j+1}(1-q_i)^{N-j}\ge0.
\]
The interchange with the finite sum is justified by the retained
absolute summability. Proposition~\ref{prop:finite-equivalences} now
gives $H_{00}$. This is the source's retained RH-to-zero-budget
construction, exposed to both routes as \code{H00\_of\_nativeRH}, and
supplies the reverse implication of Theorem~\ref{thm:main}. The remainder of the paper gives a
different forward proof. It uses Proposition~\ref{prop:finite-equivalences}
but does not invoke~\eqref{eq:retained-spectral} or the first forward
RH endpoint.

\section{From positive Bernstein rows to a probability measure}\label{sec:measure}

Write $I=[0,1]$.  For the moment, let $u=(u_n)_{n\geq0}$ be an arbitrary
real sequence and put
\[
 w^u_{N,j}=\binom Nj\sum_{r=0}^{N-j}(-1)^r\binom{N-j}{r}u_{j+r},
 \qquad 0\leq j\leq N.
\]
Define $q_{N,j}=j/N$ when $N>0$, and define $q_{0,0}=0$.
The finite Bernstein average is
\begin{equation}\label{eq:finite-average}
 A_N^u(f)=\sum_{j=0}^{N}w^u_{N,j}f(q_{N,j}).
\end{equation}
In particular, $A_0^u(f)=u_0f(0)$.  This degree-zero convention has no effect on limits as $N\to\infty$.

\begin{lemma}[Finite moment convergence]\label{lem:finite-moment-limit}
For every real sequence $u$ and every integer $k\geq0$,
\[
 \lim_{N\to\infty} A_N^u(q\mapsto q^k)=u_k.
\]
No positivity or representing-measure assumption is required.
\end{lemma}
\begin{proof}
Let $u^+_j=u_{j+1}$.  The elementary binomial identity
$(j+1)\binom{N+1}{j+1}=(N+1)\binom Nj$ gives the exact shifted-row relation
\[
 (j+1)w^u_{N+1,j+1}=(N+1)w^{u^+}_{N,j}.
\]
The term at $j=0$ contributes zero to a positive-degree monomial.  Hence
\[
 A_{N+1}^u(q^{k+1})=
 \sum_{j=0}^Nw^{u^+}_{N,j}\left(\frac{j+1}{N+1}\right)^k.
\]
For $N\geq1$, expand
\[
 \frac{j+1}{N+1}=\frac{N}{N+1}\frac jN+\frac1{N+1}
\]
by the binomial theorem.  Interchanging two finite sums yields
\begin{equation}\label{eq:moment-transition}
 A_{N+1}^u(q^{k+1})=
 \sum_{i=0}^k T_{N,k,i} A_N^{u^+}(q^i),\qquad
 T_{N,k,i}=\binom ki\left(\frac{N}{N+1}\right)^i
                   \left(\frac1{N+1}\right)^{k-i}.
\end{equation}
For $0\leq i\leq k$, $T_{N,k,i}\to\mathbf{1}_{i=k}$.
The row-mass identity $\sum_jw^u_{N,j}=u_0$ proves the degree-zero case.
Strong induction on the degree, applied to the shifted sequence $u^+$
in \eqref{eq:moment-transition}, now gives the required limit.
\end{proof}

Assume henceforth that $u_0=1$ and that every Bernstein weight is
nonnegative.  The row-mass identity makes each $A_N^u$ a positive,
unital linear functional on $C(I,\mathbb R)$, and
\begin{equation}\label{eq:average-bound}
 |A_N^u(f)|\leq\|f\|_\infty.
\end{equation}
Only values on the finite grid enter these averages.

\begin{proposition}[A constructed Hausdorff measure]\label{prop:hausdorff}
If $u_0=1$ and $w^u_{N,j}\geq0$ for all $0\leq j\leq N$, there is a
Borel probability measure $\sigma$ on $I$ such that
\begin{equation}\label{eq:hausdorff-moments}
 \int_I q^k\,d\sigma(q)=u_k\qquad(k\geq0).
\end{equation}
\end{proposition}
\begin{proof}
Fix one free ultrafilter $\mathcal U$ on $\mathbb N$ refining the
cofinite filter.  By \eqref{eq:average-bound}, the sequence $A_N^u(f)$
lies in the compact interval $[-\|f\|_\infty,\|f\|_\infty]$.  Its
$\mathcal U$-limit therefore exists for every $f\in C(I,\mathbb R)$.
Set
\[
 L(f)=\lim_{N\to\mathcal U}A_N^u(f).
\]
The same ultrafilter is used for every $f$.  Uniqueness of limits and the
finite linear identities imply linearity of $L$.  Order passes to the
limit, so $L$ is positive; moreover, $L(1)=1$ and
$|L(f)|\leq\|f\|_\infty$.  The Riesz--Markov--Kakutani representation
theorem on the compact space $I$ supplies a positive measure $\sigma$
with $L(f)=\int_I f\,d\sigma$.  Taking $f=1$ gives $\sigma(I)=1$.
Finally, ordinary convergence in Lemma~\ref{lem:finite-moment-limit}
implies convergence to the same value along $\mathcal U$, giving
\eqref{eq:hausdorff-moments}.
\end{proof}

For the actual normalized moment sequence, $u_0=1$ is a retained
unconditional identity and $H_{00}$ implies all the required Bernstein
signs.  Thus Proposition~\ref{prop:hausdorff} constructs the measure
from $H_{00}$; measure existence is not an extra premise.  This classical
construction uses the fixed free ultrafilter and the Riesz representation
theorem, and should not be interpreted as a terminating finite checker.

\section{The analytic resolvent and its actual germ}

Let $\kappa>0$ be the actual normalizing constant and let $\sigma$ be
the probability measure just constructed.  Define
\begin{equation}\label{eq:resolvent}
 S(z)=\kappa\int_I\frac{d\sigma(q)}{1+\kappa qz},\qquad
 D_\kappa=\{z\in\mathbb C:\ 1+\kappa qz\ne0\ \text{for all }q\in I\}.
\end{equation}
Certainly $0\in D_\kappa$.  Also $\mathbb H=\{z:\operatorname{Im}z>0\}$
 is contained in $D_\kappa$: for $q>0$ the denominator has positive
imaginary part, while at $q=0$ it is $1$.

\begin{lemma}[Kernel differentiation]\label{lem:resolvent-derivatives}
The function $S$ is analytic on $D_\kappa$, and for every $n\geq0$,
\begin{equation}\label{eq:resolvent-derivatives}
 S^{(n)}(z)=(-1)^n n!\,\kappa^{n+1}
 \int_I\frac{q^n}{(1+\kappa qz)^{n+1}}\,d\sigma(q)
 \qquad(z\in D_\kappa).
\end{equation}
\end{lemma}
\begin{proof}
Fix $z_0\in D_\kappa$.  Compactness of $I$ gives
$a=\min_{q\in I}|1+\kappa qz_0|>0$.  If
$|z-z_0|<a/(2\kappa)$, the reverse triangle inequality gives
$|1+\kappa qz|\geq a/2$ uniformly in $q$.  This also proves that
$D_\kappa$ is open.  For
\[
 K_n(z,q)=\frac{q^n}{(1+\kappa qz)^{n+1}},
\]
one has $\partial_z K_n=-(n+1)\kappa K_{n+1}$ and, on the indicated
neighborhood, $|K_n(z,q)|\leq(a/2)^{-n-1}$.  The derivative is likewise
bounded by a constant independent of $q$.  Since $\sigma$ is finite,
these constants are integrable majorants.  Differentiation under the
integral is therefore valid.  It proves complex differentiability locally
throughout $D_\kappa$, hence analyticity, and induction yields
\eqref{eq:resolvent-derivatives}.
\end{proof}

For $z\in\mathbb H$, the elementary kernel identity gives
\begin{equation}\label{eq:resolvent-sign}
 \operatorname{Im}S(z)
 =-\kappa^2\operatorname{Im}z
   \int_I\frac{q}{|1+\kappa qz|^2}\,d\sigma(q)\leq0.
\end{equation}
Positivity of the measure is essential here.

Let $R=\Theta'/\Theta$ be the actual logarithmic derivative.  The
retained nonvanishing theorem $\Theta(0)\ne0$ makes $R$ analytic near
zero.  The definition of the actual normalized moments, together with
the retained reality theorem for the complex moments, gives
\[
 R^{(n)}(0)=(-1)^n n!\,\kappa^{n+1}u_n.
\]
On the other hand, evaluating \eqref{eq:resolvent-derivatives} at zero
and using \eqref{eq:hausdorff-moments} gives the same expression for
$S^{(n)}(0)$.  Both functions are analytic there; their convergent Taylor
expansions consequently agree on a neighborhood of zero:
\begin{equation}\label{eq:actual-germ}
 R(z)=S(z)\quad\text{for all sufficiently small }|z|.
\end{equation}
This is equality of analytic germs, not merely equality of formal
coefficient lists.  No global zero-free assertion about $\Theta$ was
used.  The germ \eqref{eq:actual-germ}, analyticity on $\mathbb H$, and
the sign \eqref{eq:resolvent-sign} are the inputs to the analytic-defect
continuation and pole-exclusion argument.

\section{Continuation and exclusion of nonreal zeros}\label{sec:closure}
\subsection{Continuation of the analytic defect}
Write $\mathbb H=\{z\in\mathbb C:\operatorname{Im}z>0\}$. The measure construction provides a function $S$ holomorphic on $\mathbb H$ and on a neighborhood of the origin, with
\[
 \operatorname{Im}S(z)\leq 0\quad(z\in\mathbb H),
 \qquad S(z)=R(z)\quad\hbox{for all sufficiently small }z,
\]
where $R=\Theta'/\Theta$ is used only at points where $\Theta\ne0$. The latter germ is well-defined because the retained actual-function development proves $\Theta(0)\ne0$. We now extend the germ in a form that is holomorphic even at possible zeros of $\Theta$.

\begin{lemma}[Analytic-defect continuation]
Let $F$ be entire and let $S$ be holomorphic on $\mathbb H$. If there exists $\varepsilon>0$ such that $F'=SF$ on $\mathbb H\cap B(0,\varepsilon)$, then $F'=SF$ throughout $\mathbb H$.
\end{lemma}
\begin{proof}
The function $D=F'-SF$ is holomorphic on $\mathbb H$. Choose $\varepsilon>0$ as in the hypothesis, so $D=0$ on $\mathbb H\cap B(0,\varepsilon)$. Although $0$ lies on the boundary of $\mathbb H$, the point $z_*=i\varepsilon/2$ lies in both $\mathbb H$ and $B(0,\varepsilon)$. Hence $D$ vanishes on a neighborhood of $z_*$ contained in $\mathbb H$. The upper half-plane is convex and therefore connected. The identity theorem gives $D\equiv0$ on $\mathbb H$.
\end{proof}

Apply the lemma with $F=\Theta$. Continuity and $\Theta(0)\ne0$ give an origin neighborhood on which $\Theta$ does not vanish. The equality of germs $R=S$ therefore implies $\Theta'=S\Theta$ there, and analytic-defect continuation yields
\begin{equation}
 \Theta'(z)=S(z)\Theta(z)\qquad(z\in\mathbb H).
 \label{eq:theta-ode}
\end{equation}
In particular,
\begin{equation}
 \frac{\Theta'(z)}{\Theta(z)}=S(z),\qquad
 \operatorname{Im}\frac{\Theta'(z)}{\Theta(z)}\leq0
 \quad(z\in\mathbb H,\ \Theta(z)\ne0).
 \label{eq:logderiv-sign}
\end{equation}
The continuation argument requires no assertion that $\Theta$ is already zero-free: it is applied to $D$, not to a quotient across its possible poles.

\subsection{The sign condition excludes nonreal zeros}
\begin{lemma}[Upper-half-plane pole obstruction]
Let $F$ be a nonzero entire function. Suppose
\[
 \operatorname{Im}\frac{F'(z)}{F(z)}\leq0
 \qquad(z\in\mathbb H,\ F(z)\ne0).
\]
Then $F$ has no zero in $\mathbb H$.
\end{lemma}
\begin{proof}
Assume that $F(\rho)=0$ with $\operatorname{Im}\rho>0$. Since $F$ is entire and not identically zero, its zero at $\rho$ has a finite multiplicity $m\geq1$. On a neighborhood of $\rho$ we have
\[
 F(z)=(z-\rho)^m g(z),\qquad g(\rho)\ne0,
\]
with $g$ holomorphic. Shrinking the neighborhood makes $g$ nonvanishing there. Thus, for $z\ne\rho$ in that neighborhood,
\begin{equation}
 \frac{F'(z)}{F(z)}=\frac{m}{z-\rho}+\frac{g'(z)}{g(z)}.
 \label{eq:local-logarithmic-pole}
\end{equation}
Set $B=g'/g$. Continuity provides $K>0$ and $\varepsilon_0>0$, with $\varepsilon_0<\operatorname{Im}\rho$, such that
\[
 \bigl|\operatorname{Im}B(\rho-i\varepsilon)\bigr|\leq K
 \qquad(0<\varepsilon<\varepsilon_0).
\]
The same choice ensures that $F(\rho-i\varepsilon)\ne0$. Choose
\[
 0<\varepsilon<\min\{\varepsilon_0,m/K\}.
\]
Then $z_\varepsilon=\rho-i\varepsilon$ remains in $\mathbb H$, while \eqref{eq:local-logarithmic-pole} gives
\[
 \operatorname{Im}\frac{F'(z_\varepsilon)}{F(z_\varepsilon)}
 =\frac{m}{\varepsilon}+\operatorname{Im}B(z_\varepsilon)
 \geq\frac{m}{\varepsilon}-K>0.
\]
This contradicts the assumed sign condition.
\end{proof}

For the actual function $\Theta$, nontriviality follows from $\Theta(0)\ne0$. Applying the lemma to \eqref{eq:logderiv-sign} excludes all upper-half-plane zeros. The retained conjugate-zero symmetry
\[
 \Theta(z)=0\quad\Longrightarrow\quad\Theta(\overline z)=0
\]
then excludes lower-half-plane zeros as well. Consequently every zero of $\Theta$ is real. Only symmetry of the zero set is needed at this step; a separate global Schwarz-reflection hypothesis is not added to the endpoint.

\subsection{Return to the Riemann hypothesis}
The actual-function development supplies the centered-square identity
\[
 \Theta\bigl((s-\tfrac12)^2\bigr)=\xi(s)
\]
and the unconditional bound $|\operatorname{Im}s|>1$ whenever $\xi(s)=0$. Let $s=\sigma+it$ be such a zero. The centered-square identity and the reality of every $\Theta$-zero imply
\[
 0=\operatorname{Im}\bigl((s-\tfrac12)^2\bigr)
   =2(\sigma-\tfrac12)t.
\]
The height bound implies $t\ne0$, so $\sigma=1/2$. The retained equivalence between this statement about the actual $\xi$-zeros and Mathlib's \texttt{RiemannHypothesis} completes the second implication.

All analytic hypotheses used above are supplied for the actual objects by the preceding construction and the retained actual-function lemmas. In particular, no local-pole existence assumption, global zero-free assumption, or additional moment-identification premise remains in the final theorem. The forward proof does not invoke the first route's RH conclusion, although the two developments share the definition of $H_{00}$ and the finite positivity lemmas. Both equivalences use the same retained converse $\mathrm{RH}\Rightarrow H_{00}$.

\begin{proof}[Completion of Theorem~\ref{thm:main}]
Starting from $H_{00}$, Proposition~\ref{prop:finite-equivalences} supplies positive normalized Bernstein rows. Proposition~\ref{prop:hausdorff} constructs their probability measure; Lemma~\ref{lem:resolvent-derivatives} and equation~\eqref{eq:actual-germ} identify the actual resolvent germ. The continuation and pole arguments above give $\RH$. Together with the spectral forward proof and the common converse in Section~\ref{sec:spectral}, this proves the claimed equivalences.
\end{proof}

\section{Formalization and remaining obligations}\label{sec:formal}
The accompanying development uses Lean~4~\cite{lean4} and
Mathlib~\cite{mathlib}. Its two forward endpoint types are exactly
$H00\to\code{RiemannHypothesis}$; both equivalence declarations have
type $H00\leftrightarrow\code{RiemannHypothesis}$. Appendix~\ref{app:source}
records their precise names and the principal source correspondences.

The forward constructions have shared foundations. In particular, the
second route imports the definition of $H00$ and the direct finite
positivity lemmas from the first route's module. Its forward proof does
not call the first route's RH conclusion. The assertion of two routes
refers to these two forward mathematical constructions, not to disjoint
import trees.

\subsection{Verification record}
The frozen source package is version 0.2.3. The supplied native Windows
Lean~4.33.0 console excerpt reports
\code{PROJECT\_BUILD=PASS} and \code{RUN=PASS}, with three compiled
tasks and 115 reused tasks. The preserved plan contains 117 project
modules and a final audit source, hence 118 Lean files. The final audit
selects 641 declarations, including all four endpoints.

This native result is recorded as a \emph{reported} compiler-and-audit
pass: the original complete native logs and the individual axiom lists
were not independently supplied with that checkpoint. Inspection of the
preserved audit plan and runner confirms the intended coverage and the
allowed axiom set
\[
 \{\code{propext},\ \code{Classical.choice},\ \code{Quot.sound}\}.
\]
The source is preserved without changing its theorem statements or proof
bodies. Its historical preparation comments predate the later native
PASS record.

GitHub CI is a separate verification target. It pins Lean~4.33.0 and
Mathlib commit
\begin{center}\small\ttfamily
db584cd6d46c92f209a44c0f1c829460d327499d
\end{center}
and preserves the transitive dependency lock. The native Mathlib
revision was not supplied, so equality of native and CI dependency
revisions is not assumed. The workflow verifies source hashes, builds
the preserved modules, checks the exact endpoint types, and reruns the
complete axiom audit. Missing or duplicated declarations and axioms
outside the stated set cause failure.

On 28 October 2026, GitHub run 40712534224 completed successfully at
source revision \code{c6b95fd9e3d792270026da083bddddbc72bacb22}.
The job result and full log were independently read: the project build,
all four endpoint type checks, and the complete 641-declaration axiom
audit passed. Each of the four principal endpoints reported exactly
\code{propext}, \code{Classical.choice}, and \code{Quot.sound}.
The final source-identity check also passed for all 118 preserved Lean
files. The run and retained evidence are linked in~\cite{cirun}.
This observed CI result is separate from the earlier reported native
PASS. No local Lean or Lake execution was used to prepare this manuscript.

\subsection{What is not proved}
The condition $H_{00}$ is still open as an unconditional assertion about
the actual sequence. Establishing finitely many rows does not prove its
universal quantifier. A separate propagation theorem or a sound
terminating certificate checker would be required to discharge it by
finite data.

In particular, for a proposed type of finite raw witnesses and a
terminating Boolean checker, the desired statement
\[
 \bigl(\exists W:\mathrm{FiniteWitness},\ 
                 \mathrm{check}(W)=\mathrm{true}\bigr)
       \quad\Longleftrightarrow\quad\RH
\]
has not been obtained from the constructions in this note. Soundness
must connect acceptance to the actual mathematical condition, and
completeness is needed for the reverse existence implication. No
accepted finite witness is exhibited here.

The second route presented above is the measure--resolvent variant.
The original longer Widder/angular-PDE construction is a distinct
unfinished formalization. The analytic germ and continuation arguments
proved here do not certify that separate chain.

\appendix
\section{Formal statements and source correspondence}\label{app:source}
All four principal declarations belong to
\code{RH\_HilbertLimit.TwoRoutes}. Their short names and types are:
\begin{quote}\small\ttfamily
nativeRH\_of\_H00\_route\_one :\\
\hspace*{1em}H00 $\to$ \_root\_.RiemannHypothesis\\[4pt]
H00\_iff\_nativeRH\_route\_one :\\
\hspace*{1em}H00 $\leftrightarrow$ \_root\_.RiemannHypothesis\\[4pt]
nativeRH\_of\_H00\_route\_two :\\
\hspace*{1em}H00 $\to$ \_root\_.RiemannHypothesis\\[4pt]
H00\_iff\_nativeRH\_route\_two :\\
\hspace*{1em}H00 $\leftrightarrow$ \_root\_.RiemannHypothesis
\end{quote}
The defining equation is
\[
 \code{H00 := zeroBudgetCandidate.Admissible}.
\]
The spectral characterization is
\begin{center}\small\ttfamily
canonical\_actual\_subexponential\_iff\_nativeRH,
\end{center}
and the ordinary root-limit theorem is
\begin{center}\small\ttfamily
canonical\_actual\_root\_tendsto\_constructed.
\end{center}
Both are in \path{RH_HilbertLimit/AnalyticClosure.lean}.
The repository paths below refer to the fixed source revision
in~\cite{source}. All paths beginning with \code{TwoRoutes/} are
relative to \code{RH\_HilbertLimit/}.

\begin{longtable}{@{}>{\raggedright\arraybackslash}p{0.30\textwidth}>{\raggedright\arraybackslash}p{0.66\textwidth}@{}}
\toprule
Mathematical component & Source modules\\
\midrule\endhead
Actual $\xi$, centered-square identity, entire $\Theta$
 & \path{AWPIRH/XiCore.lean}; \path{AWPIRH/ThetaEntire.lean}\\[5pt]
Conjugate zeros and nonzero zero heights
 & \path{AngularMaster/ZeroGeometry.lean}; \path{AngularMaster/ZeroFreeSeed.lean}\\[5pt]
Actual moments and literal RH binding
 & \path{RH_HilbertLimit/ActualXi.lean}\\[5pt]
Cancellation condition and spectral endpoint
 & \path{RH_HilbertLimit/ActualCancellationWitness.lean};
   \path{RH_HilbertLimit/ZeroBudgetCandidate.lean};
   \path{TwoRoutes/RouteOne.lean}\\[5pt]
Finite monomial convergence
 & \path{TwoRoutes/BernsteinMomentLimit.lean}\\[5pt]
Positive functional and exact Hausdorff moments
 & \path{TwoRoutes/HausdorffMeasure.lean}\\[5pt]
Integral differentiation, sign, actual germ
 & \path{TwoRoutes/ResolventAnalytic.lean};
   \path{TwoRoutes/MeasureResolvent.lean}\\[5pt]
Analytic-defect continuation and actual pole closure
 & \path{TwoRoutes/ActualResolventContinuation.lean};
   \path{TwoRoutes/ActualHerglotzClosure.lean};
   \path{TwoRoutes/RouteTwo.lean}\\[5pt]
Full declaration audit
 & \path{TwoRoutesAudit.lean}; \path{audit-targets.tsv}\\
\bottomrule
\end{longtable}

The frozen chronological project anchor is
\code{ZETA\_HYPERBOLIC\_ANCHOR\_5}. The native package
\code{RH\_H00\_TWO\_ROUTES\_LEAN\_v0\_2\_3\_construction.zip}
has SHA-256:
\begin{center}\small\ttfamily
578cc8ad072c79dc7838fcf65a27b9eed\\
707fcc98e0fdf6a4b3601aa20e8c51a
\end{center}
The line break in this display is typographical; the digest is the
concatenation of both lines. The native checkpoint is retained as
\path{evidence/NATIVE_PASS_CHECKPOINT.json}.

\begin{thebibliography}{9}\raggedright
\bibitem{source}
Hilber-Vector-Resolvent-Theory project,
\emph{Two-route Lean source development}, version 0.2.3.
Repository revision
\code{c6b95fd9e3d792270026da083bddddbc72bacb22}, 2026.
\url{https://github.com/randogelist/Hilber-Vector-Resolvent-Theory/tree/c6b95fd9e3d792270026da083bddddbc72bacb22}.

\bibitem{lean4}
L.~de Moura and S.~Ullrich,
The Lean 4 theorem prover and programming language.
In \emph{Automated Deduction---CADE 28},
Lecture Notes in Computer Science 12699, Springer, 2021, pp.~625--635.
\url{https://doi.org/10.1007/978-3-030-79876-5_37}.

\bibitem{mathlib}
The mathlib Community,
The Lean mathematical library.
In \emph{Proceedings of the 9th ACM SIGPLAN International Conference on
Certified Programs and Proofs (CPP 2020)}, ACM, 2020, pp.~367--381.
\url{https://doi.org/10.1145/3372885.3373824}.

\bibitem{cirun}
Hilber-Vector-Resolvent-Theory project,
\emph{GitHub Actions: Lean build and axiom audit},
run 40712534224, source revision as in~\cite{source}.
\url{https://github.com/randogelist/Hilber-Vector-Resolvent-Theory/actions/runs/40712534224}.
\end{thebibliography}
\end{document}
